% Auto-generated by scripts/85_output_misspecification.R
\paragraph{Why linear endogeneity is invariant.}
Suppose the severity-indexed data can be written as \(y_{\lambda}=X\beta+a_{\lambda}X\gamma+b_{\lambda}u\). For the full sample, \(\widehat\beta_{\lambda}=\beta+a_{\lambda}\gamma+b_{\lambda}(X'X)^{-1}X'u\). After deleting any fixed set \(S\), the fitted coefficient has the same additive \(a_{\lambda}\gamma\) shift, so the full-versus-deleted coefficient difference is multiplied only by \(b_{\lambda}\). 
The OLS standard error is also multiplied by \(|b_{\lambda}|\). Therefore the standardized deletion statistic satisfies
\[
\frac{\left|\widehat\beta_{\lambda}-\widehat\beta_{\lambda,-S}\right|}
{\operatorname{SE}(\widehat\beta_{\lambda})}
=
\frac{\left|\widehat\beta_{0}-\widehat\beta_{0,-S}\right|}
{\operatorname{SE}(\widehat\beta_{0})}.
\]
Thus a departure that changes only the fitted-span coefficient and an overall residual scale leaves the standardized MIS statistic invariant. Because this equality holds for every admissible deletion set \(S\), maximizing the standardized statistic over \(S\) preserves the same invariance. The common-random-number simulation checks this proposition numerically.
